\chapter*{Εκτεταμένη Ελληνική Περίληψη}
\markboth{Εκτεταμένη Ελληνική Περίληψη}{Εκτεταμένη Ελληνική Περίληψη}
\addcontentsline{toc}{chapter}{Εκτεταμένη Ελληνική Περίληψη}

\section*{Εισαγωγή}

Η ραγδαία ανάπτυξη των υπολογιστικών υποδομών cloud και η αυξανόμενη χρήση FPGAs ως επιταχυντών υπολογισμού έχουν εισάγει νέες και ουσιαστικές προκλήσεις ασφάλειας. Στο παραδοσιακό μοντέλο απειλής, οι επιθέσεις side-channel απαιτούν φυσική πρόσβαση στη συσκευή-στόχο, εξειδικευμένο εξοπλισμό μέτρησης και άμεση ηλεκτρική σύνδεση με το κύκλωμα \cite{standaert2009introduction}. Η ενσωμάτωση των FPGAs σε περιβάλλοντα πολλαπλών χρηστών αναιρεί θεμελιωδώς αυτές τις υποθέσεις. Σε ένα κοινόχρηστο FPGA, πολλαπλοί χρήστες εκτελούν εργασίες ταυτόχρονα στο ίδιο φυσικό κύκλωμα. Ενώ η χωρική απομόνωση μεταξύ ενοικιαστών επιβάλλεται μέσω pblocks, το κοινόχρηστο δίκτυο διανομής ισχύος (Power Distribution Network, PDN) παραμένει ένα αδύνατο σημείο. Η μεταγωγική δραστηριότητα ενός χρήστη προκαλεί διακυμάνσεις τάσης που διαδίδονται μέσω του PDN και είναι ανιχνεύσιμες από οποιοδήποτε άλλο στοιχείο στο ίδιο κύκλωμα \cite{ahmed2022multi,zhao2018fpga}. Ένας κακόβουλος χρήστης μπορεί να αναπτύξει έναν αισθητήρα τάσης εντός του ολοκληρωμένου κυκλώματος και να ανακτήσει κρυπτογραφικά κλειδιά από την κατανάλωση ισχύος ενός θύματος που εκτελεί κρυπτογράφηση, χωρίς καμία φυσική παρέμβαση. Η παρούσα διπλωματική εργασία αντιμετωπίζει ακριβώς αυτή την απειλή. Προτείνεται και αξιολογείται μια ελαφριά τεχνική απόκρυψης ισχύος βασισμένη σε κυκλώματα Enhanced Ring Oscillator with Latch (ERO-L), τα οποία παράγουν θόρυβο ανεξάρτητο από τα δεδομένα στο PDN, στοχεύοντας στην υποβάθμιση του λόγου σήματος προς θόρυβο που είναι διαθέσιμος στον επιτιθέμενο. Αξιολογούνται τρεις στρατηγικές ενεργοποίησης αυξανόμενης πολυπλοκότητας, η αποτελεσματικότητα των οποίων επαληθεύεται μέσω εκτεταμένων πειραμάτων σε πραγματικό υλικό, χρησιμοποιώντας ως πρωτεύουσα μετρική την εξέλιξη κατάταξης κλειδιού.

\section*{Θεωρητικό Υπόβαθρο}

\subsection*{AES}

Η κρυπτογράφηση αποτελεί θεμελιώδη μηχανισμό προστασίας της εμπιστευτικότητας δεδομένων. Το Πρότυπο Προχωρημένης Κρυπτογράφησης (AES) είναι ο πλέον διαδεδομένος αλγόριθμος συμμετρικής κρυπτογράφησης παγκοσμίως, τυποποιημένος από το NIST το 2001 ως διάδοχος του DES \cite{fips2001advanced}. Λειτουργεί σε μπλοκ 128 bits υποστηρίζοντας κλειδιά 128, 192 και 256 bits, με 10, 12 ή 14 γύρους μετασχηματισμών αντίστοιχα \cite{daemen2002design}. Κάθε γύρος κρυπτογράφησης, εκτός του τελευταίου, αποτελείται από τέσσερις πράξεις: SubBytes (μη γραμμική αντικατάσταση bytes μέσω του πίνακα S-box), ShiftRows (κυκλική ολίσθηση γραμμών), MixColumns (γραμμικός μετασχηματισμός στηλών) και AddRoundKey (συνδυασμός κατάστασης με υπο-κλειδί μέσω XOR) \cite{fips2001advanced}. Για την παρούσα εργασία κρίσιμη είναι η ακολουθία AddRoundKey–SubBytes: το αποτέλεσμα αυτής της σύνθεσης αποτελεί το ενδιάμεσο μέγεθος που εκμεταλλεύεται η επίθεση CPA, καθώς η εξαρτημένη από τα δεδομένα κατανάλωση ισχύος κατά τον υπολογισμό του S-box συσχετίζεται στατιστικά με το μυστικό κλειδί \cite{lo2017power,brier2004correlation}.

\subsection*{Αρχιτεκτονική FPGA και Ασφάλεια}

Τα FPGAs (Field Programmable Gate Arrays) είναι επαναπρογραμματιζόμενα ψηφιακά κυκλώματα που αποτελούνται από πλέγμα λογικών μπλοκ (CLBs), καθένα εκ των οποίων περιέχει πίνακες αναζήτησης (LUTs) και flip-flops, συνδεδεμένα μέσω προγραμματιζόμενου δικτύου δρομολόγησης. Η εξέλιξη των FPGA τα έχει μετατρέψει από απλές συστοιχίες λογικής σε σύνθετα συστήματα με ενσωματωμένες μνήμες RAM, μπλοκ DSP και υποσυστήματα επεξεργαστή \cite{boutros2021fpga}. 

Σύγχρονα FPGAs σε πλατφόρμες cloud επιτρέπουν κοινή χρήση μιας συσκευής από πολλαπλούς χρήστες. Αυτή η κοινή χρήση εισάγει νέες ευπάθειες ασφάλειας: σε αντίθεση με επεξεργαστές γενικής χρήσης, τα FPGAs επιτρέπουν στους χρήστες να αναπτύσσουν αυθαίρετο υλικό, ενώ μοιράζονται φυσικούς πόρους όπως το PDN. Ένας κακόβουλος χρήστης μπορεί να εκμεταλλευτεί αυτό το κοινόχρηστο μέσο για επιθέσεις side-channel, εισαγωγή σφαλμάτων ή επιθέσεις άρνησης υπηρεσίας (denial of service) \cite{ahmed2022multi}.

\subsection*{Επιθέσεις Side-Channel και CPA}

Οι επιθέσεις side-channel (SCAs) εκμεταλλεύονται φυσικές πληροφορίες που διαρρέουν από την υλοποίηση ενός συστήματος, αντί να επιτίθενται στον ίδιο τον κρυπτογραφικό αλγόριθμο \cite{gattu2020power}. Τα κανάλια ισχύος εκμεταλλεύονται παραλλαγές κατανάλωσης ισχύος κατά τη διάρκεια υπολογισμών, τα χρονικά κανάλια παραλλαγές χρόνου εκτέλεσης και τα ηλεκτρομαγνητικά κανάλια εκπομπές από ηλεκτρονικά εξαρτήματα \cite{gandolfi2001electromagnetic}. Η ανάπτυξη FPGAs σε υποδομές ψλοθδεισάγει νέες ευπάθειες που επεκτείνουν τις SCAs πέρα από παραδοσιακές πλατφόρμες ενσωματωμένων συστημάτων \cite{zhao2018fpga}.

Η Correlation Power Analysis (CPA) εισήχθη επίσημα από τους Brier, Clavier και Olivier \cite{brier2004correlation} και βασίζεται στο μοντέλο Hamming Weight: η κατανάλωση ισχύος θεωρείται ανάλογη του αριθμού των bits που τίθενται σε 1 κατά τη διάρκεια μιας πράξης \cite{lo2017power}. Η επίθεση εκτελείται σε τέσσερις φάσεις: συλλογή ιχνών ισχύος ενώ επεξεργάζονται γνωστές είσοδοι, κατασκευή υποθέσεων Hamming Weight για κάθε πιθανό υπό-κλειδί, υπολογισμός συντελεστή συσχέτισης Pearson μεταξύ προβλεπόμενων και μετρούμενων τιμών, και ανάκτηση κλειδιού ως εκείνου με τη μέγιστη απόλυτη συσχέτιση \cite{lo2017power,brier2004correlation}.

% \begin{figure}[ht]
%     \centering
%     \includegraphics[width=1.0\textwidth]{Figures_2/attack_pipeline.png}
%     \caption*{Σύνοψη της μεθοδολογίας της CPA}
% \end{figure}

Κρίσιμο πλεονέκτημα της CPA είναι ότι επιτρέπει ανεξάρτητη επίθεση σε κάθε ένα από τα 16 bytes του AES-128, ανάγοντας το πρόβλημα αναζήτησης $2^{128}$ πιθανοτήτων σε 16 ανεξάρτητες αναζητήσεις 256 πιθανοτήτων\cite{lo2017power}. Για την αξιολόγηση αντιμέτρων χρησιμοποιείται η μετρική κατάταξης κλειδιού (key rank), που εκτιμά τη θέση του σωστού κλειδιού στη λίστα όλων των πιθανών κλειδιών ταξινομημένων κατά πιθανότητα βάσει αποτελεσμάτων CPA \cite{spielmann2023rds}. Τιμή $2^{128}$ σημαίνει μηδενική πληροφορία για τον επιτιθέμενο, τιμή 0 σημαίνει πλήρης ανάκτηση.

\section*{Μεθοδολογία}

\subsection*{Επιλογή κυκλώματος Ring-Oscillator} 
Η επιλογή κατάλληλου κυκλώματος για τη απόκρυψη ισχύος απαιτούσε δομημένη αξιολόγηση. Πέντε υποψήφια κυκλώματα ταλαντωτή δακτυλίου, ERO-cf, ERO-L-cf, ERO-L, ERO-FF-cf-L και ERO-FF-L, αξιολογήθηκαν ως προς πέντε κριτήρια: ένταση μεταγωγής, αξιοποίηση δρομολόγησης και φόρτωση, σταθερότητα ταλάντωσης, συμμόρφωση με κανόνες σύνθεσης και συμβατότητα με λειτουργικά συστήματα \cite{mahmoud2025roboost}.

Κάθε κύκλωμα υλοποιήθηκε μεμονωμένα στο FPGA και η συμπεριφορά του καταγράφηκε μέσω του αισθητήρα RDS, επιτρέποντας άμεση οπτική σύγκριση με τη βασική γραμμή AES. Τα ERO-cf, ERO-L-cf και ERO-FF-cf-L περιέχουν άμεσο συνδυαστικό βρόχο ανατροφοδότησης, προκαλώντας κρίσιμη παραβίαση DRC τύπου LUTLP-1 στο Vivado. Αυτό τα καθιστά ακατάλληλα για ανάπτυξη σε περιβάλλοντα με αυστηρή επιβολή κανόνων σύνθεσης. Μεταξύ των συμβατών επιλογών, το ERO-L υπερτερεί σαφώς έναντι του ERO-FF-L: παρουσιάζει μέση έξοδο αισθητήρα 108-109 μονάδων, έναντι ~103 μονάδων για το ERO-FF-L, λόγω χρήσης latch αντί flip-flop στο μονοπάτι ανατροφοδότησης. Ο latch λειτουργεί σε διαφανή κατάσταση, διαδίδοντας αλλαγές ταχύτερα και διατηρώντας υψηλότερη συχνότητα ταλάντωσης \cite{mahmoud2025roboost}. 

Το κύκλωμα ERO-L, αποτελείται από τέσσερα LUTs (A, B, C, D) ρυθμισμένα ως πύλες NAND και διασυνδεδεμένα σε πλήρως συνδεδεμένη τοπολογία με υψηλό fanout. Ένας latch εισάγεται στο μονοπάτι ανατροφοδότησης από τη LUT A, με εξωτερικό σήμα ενεργοποίησης που επιτρέπει άμεσο έλεγχο της ταλάντωσης. Αυτό το σχέδιο επιτυγχάνει υψηλή ένταση μεταγωγής ενώ παράλληλα αποφεύγει τον άμεσο συνδυαστικό βρόχο, εξασφαλίζοντας πλήρη συμβατότητα με τα εργαλεία σύνθεσης.

\subsection*{Στρατηγικές Ενεργοποίησης} 

Αφού επιλέχθηκε το ERO-L, το επόμενο ερώτημα ήταν πώς και πότε πρέπει να ενεργοποιείται σε σχέση με την κρυπτογράφηση AES που προστατεύει. Ο στόχος είναι να αλλοιωθούν τα ίχνη που καταγράφει ο αισθητήρας RDS του αντιπάλου. Η στρατηγική ενεργοποίησης καθορίζει τη χρονική δομή του παραγόμενου θορύβου και άρα πόσο αποτελεσματικά μπορεί ο θόρυβος να αντισταθεί στη διαδικασία συσχέτισης της CPA \cite{mangard2007power}. Αξιολογήθηκαν τρεις στρατηγικές αυξανόμενης πολυπλοκότητας:

\textbf{Στατική Ενεργοποίηση:} Το σήμα ενεργοποίησης είναι μόνιμα ενεργό καθ' όλη τη διάρκεια του πειράματος. Είναι η απλούστερη υλοποίηση, χωρίς λογική ελέγχου χρονισμού. Επειδή το ERO-L είναι ήδη ενεργό κατά τη βαθμονόμηση, η βαθμονόμηση γίνεται υπό τις πραγματικές συνθήκες λειτουργίας. Η ουσιαστική αδυναμία είναι ότι η σταθερή συνεισφορά του ERO-L αφαιρείται αυτόματα κατά τη φάση αφαίρεσης μέσης τιμής της CPA, αφήνοντας μόνο ένα αυξημένο επίπεδο θορύβου ως εμπόδιο για τον αντίπαλο \cite{mangard2007power}. Αυτό σημαίνει ότι η στατική στρατηγική δεν εισάγει ουσιαστική ίχνος-προς-ίχνος παραλλαγή, που είναι απαραίτητη για να αντισταθεί σε CPA.

\textbf{Περιοδική Ενεργοποίηση:} Το σήμα ενεργοποίησης εναλλάσσεται on/off με σταθερή συχνότητα και κύκλο λειτουργίας 50\% κατά τη διάρκεια κάθε ίχνους. Υλοποιείται μέσω αποκλειστικής λογικής ελέγχου συγχρονισμένης με το ρολόι του αισθητήρα, παραμετροποιημένης ως προς αριθμό δειγμάτων ανά ίχνος, αριθμό περιόδων και κύκλο λειτουργίας. Το σήμα ενεργοποίησης μεταφέρεται στον τομέα κύριου ρολογιού μέσω συγχρονιστή δύο flip-flop. Η στρατηγική εισάγει χρονική παραλλαγή εντός κάθε ίχνους, αλλά το μοτίβο είναι ταυτόσημο σε όλα τα ίχνη. Αυτό αποτελεί θεμελιώδη αδυναμία, καθώς ένας πληροφορημένος αντίπαλος θα μπορούσε θεωρητικά να ανακτήσει και να αφαιρέσει το περιοδικό μοτίβο. Επιπλέον, αν η επιλεγμένη περίοδος συνάδει με τη δομή γύρων AES, το αποτέλεσμα μπορεί να είναι αντίθετο από το επιθυμητό.

\textbf{Τυχαιοποιημένη Ενεργοποίηση:} Παράγεται μοναδική ψευδο-τυχαία ακολουθία ενεργοποίησης για κάθε ίχνος ανεξάρτητα. Ο αριθμός ενεργών παραθύρων, οι θέσεις τους και τα μήκη τους μεταβάλλονται από ίχνος σε ίχνος, εξασφαλίζοντας ότι η συνεισφορά ERO-L δεν έχει συνεκτική χρονική δομή. Η λογική ελέγχου, υλοποιείται ως μηχανή πεπερασμένων καταστάσεων τεσσάρων καταστάσεων (IDLE, DECIDE\_NUM\_WINDOWS, GENERATE\_WINDOWS, COLLECTING) οδηγούμενη από 16-bit LFSR μέγιστου μήκους. Ο LFSR τρέχει συνεχώς. Η τιμή που παρουσιάζει κατά την έναρξη κάθε ίχνους εξαρτάται από τον χρονισμό του σήματος data-ready του AES, παρέχοντας αποτελεσματικά διαφορετική αρχική τιμή ανά ίχνος. Επειδή το μοτίβο ενεργοποίησης είναι διαφορετικό για κάθε ίχνος, η συνεισφορά ERO-L δεν συσσωρεύεται συνεκτικά υπό τη διαδικασία συσχέτισης CPA. Αυτή είναι η θεμελιώδης ιδιότητα που καθιστά τη στρατηγική αυτή ανώτερη.

\subsection*{Πειραματική Υποδομή και Ροή Αξιολόγησης} 
Η υλοποίηση έγινε στην πλακέτα Digilent Nexys A7-100T που βασίζεται στο FPGA AMD Artix-7 XC7A100T (63400 LUTs, 126800 flip-flops, 6 CMTs με PLLs). Η αρχιτεκτονική συστήματος, εφαρμόζει φυσική απομόνωση μέσω περιορισμών Pblock: θύμα (AES-128), αντίπαλος (RDS[24], βαθμονόμηση, FIFO) και ERO-L τοποθετούνται σε διακριτές μη επικαλυπτόμενες περιοχές \cite{spielmann2023rds}. Το ERO-L βρίσκεται χωρικά μακριά και από τις δύο άλλες περιοχές, ώστε η επίδρασή του να διαδίδεται αποκλειστικά μέσω του κοινόχρηστου PDN. Λογική απομόνωση εξασφαλίζεται από την απουσία οποιωνδήποτε άμεσων συνδέσεων σήματος μεταξύ των τριών περιοχών.

\begin{figure}[ht]
    \centering
    \includegraphics[width=0.9\textwidth]{Figures_3/full_system_architecture.png}
    \caption*{Αρχιτεκτονική του συστήματος}
\end{figure}

Το σύστημα λειτουργεί σε δύο ανεξάρτητους τομείς ρολογιού: αισθητήρας στα 200 MHz και AES κυρίως στα 100 MHz, με πρόσθετα πειράματα στα 50 MHz και 20 MHz. Το ERO-L είναι ασύγχρονο ως προς αμφότερους τους τομείς και η συχνότητα ταλάντωσής του καθορίζεται από τις εσωτερικές του καθυστερήσεις διάδοσης, όχι από εξωτερικό ρολόι. Η ανάλυση CPA εκτελείται εκτός επεξεργαστή στον κεντρικό υπολογιστή με επιτάχυνση GPU, χρησιμοποιώντας τον αλγόριθμο εκτίμησης κατάταξης κλειδιού των Glowacz et al \cite{glowacz2015simpler}.


\section*{Πειραματικά Αποτελέσματα} 

\subsection*{Απροστάτευτη Υλοποίηση} 

Το πρώτο πείραμα αξιολόγησε την CPA επίθεση στο απροστάτευτο AES-128 στα 100 MHz. Η εξέλιξη του key rank, αποδεικνύει γρήγορη σύγκλιση: ξεκινώντας από τιμή 120, με 1000 ίχνη, η κατάταξη πέφτει ραγδαία μεταξύ 3.000 και 10000 ιχνών και φτάνει στο 0 περίπου στα 10000 ίχνη, σηματοδοτώντας πλήρη ανάκτηση κλειδιού. Το στενό χάσμα μεταξύ άνω και κάτω ορίου υποδηλώνει σταθερές συνθήκες μέτρησης και ισχυρή διαρροή. Αυτά τα αποτελέσματα αποτελούν τη βάση αναφοράς για όλες τις επόμενες συγκρίσεις.

\subsection*{Αποτελέσματα Στατικής Ενεργοποίησης}

Η στατική ενεργοποίηση αξιολογήθηκε με 128, 256 και 320 παράλληλα ERO-L. Η εξέλιξη του key rank αποκαλύπτει συμπεριφορά που αντίκειται στις αρχικές προσδοκίες. Τα 128 στιγμιότυπα αποτελούν την πιο αποτελεσματική στατική διαμόρφωση: η κατάταξη φτάνει περίπου στο 72 στα 10000 ίχνη χωρίς σύγκλιση. Τα 256 στιγμιότυπα παρουσιάζουν πιο σύνθετη συμπεριφορά, με ταχύτερη αρχικά αλλά βραδύτερη κατόπιν κατάβαση. Τα 320 στιγμιότυπα, η διαμόρφωση με τη μεγαλύτερη ορατή συσκότιση σε επίπεδο ίχνους, συγκλίνουν στο 0 εντός 10000 ιχνών, χειρότερα από την απροστάτευτη βασική γραμμή.

\begin{figure}[h!]
    \centering
    \includegraphics[width=0.8\textwidth]{Figures_4/keyrank_static.png}
    \caption*{Αποτελέσματα key rank για στατική ενεργοποίηση}
\end{figure}

Αυτό το αποτέλεσμα εξηγείται από τον δομικό περιορισμό της στατικής ενεργοποίησης: η σταθερή συνεισφορά ERO-L αφαιρείται από την προεπεξεργασία της CPA και η βαρύτερη φόρτωση PDN σε μεγαλύτερο αριθμό στιγμιοτύπων μετατοπίζει το σημείο λειτουργίας του αισθητήρα με τρόπους που ενδέχεται να ενισχύουν ορισμένα χαρακτηριστικά διαρροής. Το αποτέλεσμα αυτό υπογραμμίζει ότι η αύξηση αριθμού στιγμιοτύπων υπό στατική ενεργοποίηση δεν μεταφράζεται αξιόπιστα σε ισχυρότερη αντίσταση CPA.

\subsection*{Αποτελέσματα Περιοδικής Ενεργοποίησης} 

Η περιοδική ενεργοποίηση αξιολογήθηκε με τρεις διαμορφώσεις περιόδου (2, 4 και 6 πλήρεις περίοδοι ανά ίχνος 256 δειγμάτων, κύκλος λειτουργίας 50\%). Η ανάλυση εξέλιξης του key rank για τις τρεις περιόδους παρουσιάζει μη μονοτονία: η διαμόρφωση 4 περιόδων είναι η πιο επιβλαβής, καθώς όλα τα πλήθη ERO-Ls επιταχύνουν τη σύγκλιση σε σχέση με την απροστάτευτη υλοποίηση, με τα 320 να φτάνουν στο μηδέν περί τα 7000 ίχνη. Η περίοδος 64 δειγμάτων συντονίζεται δυσμενώς με τη δομή γύρων AES στα 100 MHz, με αποτέλεσμα οι μεταβάσεις on/off να εμπίπτουν με συνέπεια στις ίδιες υπολογιστικές φάσεις σε όλα τα ίχνη, ενισχύοντας τη διαρροή αντί να τη μειώνει. Πρόκειται για σαφή εφαρμογή του θεωρητικά αναμενόμενου κινδύνου.

\begin{figure}[h!]
    \centering
    \includegraphics[width=0.8\textwidth]{Figures_4/keyrank_per_2.png}
    \caption*{Αποτελέσματα key rank για περιοδική ενεργοποίηση, 128 δείγματα/περίοδο}
\end{figure}

\begin{figure}[h!]
    \centering
    \includegraphics[width=0.8\textwidth]{Figures_4/keyrank_per_4.png}
    \caption*{Αποτελέσματα key rank για περιοδική ενεργοποίηση, 64 δείγματα/περίοδο}
\end{figure}

\begin{figure}[h!]
    \centering
    \includegraphics[width=0.8\textwidth]{Figures_4/keyrank_per_6.png}
    \caption*{Αποτελέσματα key rank για περιοδική ενεργοποίηση, 42 δείγματα/περίοδο}
\end{figure}

Η διαμόρφωση 6 περιόδων είναι η πιο ευνοϊκή: τα 256 ERO-Ls διατηρούν κατάταξη άνω του 80 στα 10000 ίχνη χωρίς εμφανή τάση σύγκλισης. Η διαμόρφωση 2 περιόδων βρίσκεται ενδιάμεσα με ανάμεικτα αποτελέσματα. Το κύριο συμπέρασμα είναι ότι η σχέση μεταξύ παραμέτρων περιόδου και αμυντικής αποτελεσματικότητας δεν μπορεί να προβλεφθεί από οπτική ανάλυση ιχνών και ότι κακώς επιλεγμένη παράμετρος μπορεί να μειώσει δραστικά την ασφάλεια.

\subsection*{Αποτελέσματα Τυχαιοποιημένης Ενεργοποίησης} 
Η τυχαιοποιημένη ενεργοποίηση αξιολογήθηκε σε εκτεταμένο πείραμα 20000 ιχνών. Η εξέλιξη του key rank, επιβεβαιώνει την ισχύ της τυχαιοποιημένης ενεργοποίησης και αποκαλύπτει μια εντυπωσιακά συνεπή σχέση: όσο περισσότερα ERO-Ls, τόσο ισχυρότερη αντίσταση. Με 128 η κατάταξη φτάνει περίπου στο 25 στα 20000 ίχνη. Με 256 φτάνει περίπου στο 75 στα 20000 ίχνη χωρίς σημεία σύγκλισης. Με 320 η κατάταξη παραμένει ουσιαστικά επίπεδη κοντά στο 125 καθ' όλη τη διάρκεια του πειράματος, το ισχυρότερο αποτέλεσμα ανάμεσα σε όλες τις αξιολογούμενες στρατηγικές και αριθμούς στιγμιοτύπων.

\begin{figure}[h!]
    \centering
    \includegraphics[width=0.8\textwidth]{Figures_4/keyrank_rand.png}
    \caption*{Αποτελέσματα key rank για τυχαία ενεργοποίηση}
\end{figure}

Η βελτίωση οφείλεται στον θεμελιώδη μηχανισμό: επειδή το μοτίβο ενεργοποίησης μεταβάλλεται ανεξάρ-\\τητα για κάθε ίχνος, η συνεισφορά ERO-L δεν συσσωρεύεται συνεκτικά υπό τη διαδικασία συσχέτισης CPA ανεξαρτήτως αριθμού ERO-L. Η απουσία αποτυχιών εξαρτώμενων από ευθυγράμμιση, που κατέστησαν την περιοδική στρατηγική αναξιόπιστη, καθιστά την τυχαιοποιημένη ενεργοποίηση την πιο εύρωστη στρατηγική.

\subsection*{Επίδραση Συχνότητας Ρολογιού AES} 

Για να αξιολογηθεί η γενικευσιμότητα των αποτελεσμάτων, όλες οι στρατηγικές ενεργοποίησης δοκιμά-\\στηκαν επίσης στα 50 MHz και 20 MHz με 320 ERO-Ls. Τα αποτελέσματα αποκαλύπτουν συνεπή βελτίωση αποτελεσματικότητας σε χαμηλότερες συχνότητες για στατική και περιοδική ενεργοποίηση: η στατική ενεργοποίηση που αποτυγχάνει στα 100 MHz παρέχει ουσιαστικά πλήρη αντίσταση στα 20 MHz, και η περιοδική ενεργοποίηση ενισχύεται σε παρόμοιο βαθμό. Για τυχαιοποιημένη ενεργοποίηση η σχέση είναι λιγότερο μονότονη: καμία από τις τρεις συχνότητες δεν παρουσιάζει σύγκλιση εντός του πειραματικού προϋπολογισμού, αλλά στα 50 MHz παρατηρείται ελαφρώς βραδύτερη πτώση σε σχέση με τα 100 και 20 MHz. 

\subsection*{Συνοπτική Σύγκριση} 
Ο παρακάτω πίνακας συνοψίζει τη χρήση πόρων για τις τρεις στρατηγικές N320 στον Artix-7. Ο πίνακας ERO-L καταναλώνει 320 LUTs (0.5\%) σε όλες τις στρατηγικές. Ο ελεγκτής ενεργοποίησης μεταβάλλεται: στατική ενεργοποίηση δεν απαιτεί ελεγκτή, περιοδικός ελεγκτής προσθέτει 29 LUTs (0.05\%) και τυχαιοποιημένος ελεγκτής 129 LUTs (0.2\%). Συνολικά το σύστημα χρησιμοποιεί 3096-3225 LUTs (4.88-5.09\%), αφήνοντας άνω του 94\% των πόρων διαθέσιμο.

\begin{table}[h]
\centering
\begin{tabular}{lccc}
\hline
\textbf{Στρατηγική} & \textbf{ERO-L LUTs (\%)} & \textbf{Controller LUTs (\%)} & \textbf{Συνολικά LUTs (\%)} \\
\hline
Baseline (χωρίς ERO-L) & ---         & ---         & 2776 (4.38\%) \\
Στατική              & 320 (0.5\%) & ---         & 3096 (4.88\%) \\
Περιοδική (6-period) & 320 (0.5\%) & 29 (0.05\%) & 3125 (4.93\%) \\
Τυχαιοποιημένη          & 320 (0.5\%) & 129 (0.2\%) & 3225 (5.09\%) \\
\hline
\end{tabular}
\caption*{Απαιτούμενοι πόροι στο Artix-7 XC7A100T}
\label{tab:lut_utilization_el}
\end{table}

Ο παρακάτω πίνακας συνοψίζει την αποτελεσματικότητα της μεθόδου για N320 στα 100 MHz. Η στατική ενεργοποίηση συγκλίνει στο 0 εντός 10000 ιχνών, το ίδιο αποτέλεσμα με την απροστάτευτη βασική γραμμή, άρα αναποτελεσματική. Η περιοδική ενεργοποίηση (6 περίοδοι) διατηρεί κατάταξη ~25 στα 10000 ίχνη χωρίς σύγκλιση, ωστόσο άλλες περίοδοι είναι ενεργά επιβλαβείς. Η τυχαιοποιημένη ενεργοποίηση διατηρεί κατάταξη 120 στα 20000 ίχνη χωρίς καμία τάση σύγκλισης, η ισχυρότερη και πλέον αξιόπιστη επίδοση.

\begin{table}[h]
\centering
\begin{tabular}{lccc}
\hline
\textbf{Στρατηγική} & \textbf{Αριθμός Traces} & \textbf{Key Rank} & \textbf{Συγκλίνει?} \\
\hline
Baseline            & 10000 & 0   & Yes \\
Στατική              & 10000 & 0   & Yes \\
Περιοδική (6-period) & 10000 & 25  & No \\
Τυχαιοποιημένη          & 20000 & 120 & No \\
\hline
\end{tabular}
\caption*{Αποτελεσματικότητα των πειραμάτων}
\label{tab:obfuscation_effectiveness_el}
\end{table}
